\documentclass[11pt, a4paper, twoside]{article} % , draft
\usepackage[utf8]{inputenc}

\usepackage{enumitem} % customiçe item dots etc
\usepackage{textgreek} % obv
\usepackage{physics} % for easy derivative notation
\usepackage{amsmath}
\usepackage{amsthm} %theorems
\usepackage{amssymb}
\usepackage{mathtools} % for matrices with blocks inside
\usepackage[scr=boondoxo]{mathalfa}
\usepackage{pst-node}%
\usepackage{mathrsfs}
\DeclareMathAlphabet{\mathpzc}{OT1}{pzc}{m}{it}

\newcommand{\mc}{\multicolumn{1}{c}}
\newcommand{\R}{\mathbb{R}} % command for real R
\newcommand{\Holo}{\mathcal{H}}
\newcommand{\M}{\mathcal{M}}
\newcommand{\C}{\mathbb{C}}
\newcommand{\N}{\mathbb{N}}
\newcommand{\z}{\mathpzc{s}}
\newcommand{\p}{\mathpzc{r}}
\newcommand{\s}{\mathbb{S}}
\newcommand{\W}{\mathbb{W}}
\newcommand{\U}{\mathscr{U}}
\newcommand{\Lg}{\mathscr{L}}
\newcommand{\x}{\mathcal{X}}

\usepackage{csquotes}
\MakeOuterQuote{"}
\setlength{\parskip}{0.3 cm}

\usepackage{fancyhdr}

%\usepackage{nath} % authomatic parenthesis stuff
%\delimgrowth=1
\usepackage[left=2cm, right=2cm, top=2.1cm, bottom=2.1cm]{geometry} % set custom margins
\usepackage{graphicx} % to insert figures
\usepackage{grffile}
\graphicspath{{Figures/}} % define the figure folder path
\usepackage{subcaption} % for multiple figures at once each with a caption
\usepackage{multirow} %multirow in tables

\usepackage{caption}
\captionsetup[figure]{font=footnotesize} %adjust caption size
\captionsetup[table]{font=footnotesize} %adjust caption size

\usepackage{booktabs} % for pretty tabs in tables
\usepackage{siunitx} % Required for alignment
\captionsetup{labelfont=bf} % bold face captations

\usepackage{hyperref} % makes every reference a hyperlink
\hypersetup{
    colorlinks=true,
    linkcolor=violet,
    filecolor=[rgb]{0.69, 0.19, 0.38},      
    urlcolor=[rgb]{0.0, 0.81, 0.82},
    citecolor=[rgb]{0.69, 0.19, 0.38}
}

\usepackage{epigraph} % for quotations in teh begginig
\setlength\epigraphwidth{8cm}
\setlength\epigraphrule{0pt}
\usepackage{etoolbox}
\makeatletter
\patchcmd{\epigraph}{\@epitext{#1}}{\itshape\@epitext{#1}}{}{}
\renewcommand{\qedsymbol}{o.\textepsilon.\textdelta}

\newtheorem{prop}{Proposition} %so I can use propositions
\newtheorem{cor}{Corollary} %so I can use corollaries
\newtheorem{defi}{Definition} %so I can use corollaries

\makeatother % all this is for the epigraph
\usepackage{tocloft}

\usepackage{imakeidx} % make index

\makeindex[columns=3, title=Alphabetical Index, intoc]

%\title{\vspace{-2.5cm} {\bf Can we make the Exponential scaling in Time\\ be Linear in Time if Parallelized Exponentially? \\ {\em - Part 2 -}} \vspace{-0.4cm}  }
\title{\vspace{-2cm} {\bf End-to-End Creation of a \\Conical Refraction Polarimeter}\\{\small by {\em Xabier Oyanguren Asua}}\vspace{-0.8cm}}
\date{\vspace{-11ex}}
\let\clipbox\relax
\usepackage{adjustbox}
\newcolumntype{?}{!{\vrule width 1.5pt}}
\usepackage{abstract}
\setlength{\absleftindent}{0mm}
\setlength{\absrightindent}{0mm}

\usepackage{tcolorbox}
\DeclareRobustCommand{\mybox}[2][gray!10]{%
\begin{tcolorbox}[   %% Adjust the following parameters at will.
        left=0.2cm,
        right=0.2cm,
        top=0.15cm,
        bottom=0.15cm,
        colback=#1,
        colframe=#1,
        width=\dimexpr\textwidth\relax, 
        enlarge left by=0mm,
        boxsep=5pt,
        arc=0pt,outer arc=0pt,
        ]
        #2
\end{tcolorbox}
}


\usepackage{anyfontsize}

\NewDocumentEnvironment{kapituloBerria}{mm}
{\clearpage           % we want a new page          %% I commented this
   \thispagestyle{empty}% no header and footer
   \vspace*{\stretch{2}}% some space at the top
   \raggedleft          % flush to the right margin
   {\textbf{{\fontsize{60}{40}\selectfont \hspace{+9.5cm}#1 \newline \newline}}}
   \bf
   \fontsize{30}{20}\selectfont
  }
  {\par % end the paragraph
   \vspace{\stretch{3}} % space at bottom is three times that at the top
   \normalfont
      \fontsize{15}{20}\selectfont
      \vspace{-1cm}
   \begin{flushleft}{ \textit{#2}  }
   \end{flushleft}
   \clearpage           % finish off the page
  }
  
%\newenvironment{kapituloBerria}[2]
  

\usepackage{listings}
\usepackage{xcolor}
\lstset{language=C++,
                basicstyle=\ttfamily,
                keywordstyle=\color{blue}\ttfamily,
                stringstyle=\color{red}\ttfamily,
                commentstyle=\color{green}\ttfamily,
                morecomment=[l][\color{magenta}]{\#}
    backgroundcolor=\color{black!5}, % set backgroundcolor
    basicstyle=\footnotesize,% basic font setting
}

\begin{document}

\clearpage
%% temporary titles
% command to provide stretchy vertical space in proportion
\newcommand\nbvspace[1][3]{\vspace*{\stretch{#1}}}
% allow some slack to avoid under/overfull boxes
\newcommand\nbstretchyspace{\spaceskip0.5em plus 0.25em minus 0.25em}
% To improve spacing on titlepages
\newcommand{\nbtitlestretch}{\spaceskip0.6em}
\pagestyle{empty}
\begin{center}
\bfseries
\nbvspace[1]
\Huge
{\nbtitlestretch
{\normalsize END-TO-END CREATION OF A }\\
\Huge CONICAL REFRACTION \\
POLARIMETER }

\nbvspace[4]
{\small

BUILDING A STATE OF THE ART\vspace{0.2cm}\\ AFFORDABLE DEVICE\vspace{0.2cm}\\ TO MONITOR LINEAR POLARIZATION\\
}

\normalsize

\nbvspace[5]
\small BY\\
\Large Xabier Oyanguren Asua\\[0.5em]
%\footnotesize Can we really foresee the future of the Universe?"\\

\nbvspace[5]

\small THESIS SUPERVISOR \vspace{0.15cm}\\
\Large Àlex Turpin Avilés

\nbvspace[1]

\includegraphics[width=2.5in]{UAB.png}
\normalsize
\vspace{-0.5cm}
%\small Thesis Directors: \\
%\nbvspace[0.2]
%\large Jordi Mompart Penina \\
%Xavier Oriols Pladevall\\
%\nbvspace[1]

%Universitat Autònoma de Barcelona\\
\large
%DEGREE FINAL DISSERTATION \\
\small
Bachelor's Thesis \\ Computational Mathematics and Data Analytics \\
\vspace{0.1cm}
\small
2021-2022
\nbvspace[1]
\end{center}
\newpage
\null
\clearpage

\maketitle
\vspace{-0.3cm}
\pagenumbering{gobble}
\setlength{\cftbeforesecskip}{0.4cm}
\setlength{\cftbeforesubsecskip}{0.4cm}
\setlength{\cftbeforesubsubsecskip}{0.25cm}

\tableofcontents
\clearpage
\pagenumbering{arabic}
\setcounter{page}{1}
\vspace{-0.3 cm}

%\pagestyle{empty}
\pagestyle{fancy}

\section*{\centering \huge{Introduction}\vspace{-0.3cm}}
\noindent\rule{\textwidth}{0.4pt}
\addcontentsline{toc}{section}{Introduction}
\section*{Abstract}


\section*{Objectives}\vspace{-0.2cm}



\section*{Guideline}\vspace{-0.2cm}


\newpage

%%%%%%%%%%%%%%%%%%%%%%%%%%%%%%%%%%%%%%%%%%%%%%%%%%%%%%%%%%%%%%%%%%%%%%%%%%%%%%%%%%%%%%%%%%%%%%%%%%%%%%%%%%%%%%%%%%%%%%%%%%%%%%%%%%%%%%%%%%%%%%%%%%%%%%%%%%%%%%%%%%%%%%%%%%%%%
\newpage


%\begin{kapituloBerria}{Part A}{ Understanding the Underlying Phenomenon}
%Conical Refraction Essentials 
%\end{kapituloBerria}
%\newpage
%\fancyhead[L]{\null}
%\fancyhead[R]{\null}
%\null
%\clearpage


\section*{\centering \huge{Part A: Conical Refraction Essentials}\vspace{-0.3cm}}
\noindent\rule{\textwidth}{0.4pt}
\vspace{-0.2cm}
\addcontentsline{toc}{section}{Part A: Conical Refraction Essentials}

In this first section, we will review the basic theoretical framework predicting the Conical Refraction (CR) phenomenon on which the developed device will be based.

\section*{A.1. Derivation of Conical Refraction from Maxwell's Equations}
\subsection*{A.1.1. Electromagnetic Waves }
Given we denote by $\vec{E}(\vec{r},t)$ and $\vec{B}(\vec{r},t)$ respectively the electric and magnetic forces at the spatial position $\vec{r}=(x,y,z)\in\R^3$ and a given time $t\in\R$, their shape is given by the well known (microscopic) Maxwell Equations, ruled by the electric charge density $\vec{\rho}(\vec{r},t)$ (charge per unit volume) and electric current density $\vec{J}(\vec{r},t)$ (charge density times the velocity field of its constituent particles):
\begin{equation}\label{div}
\vec{\nabla}_r\cdot \vec{E}(\vec{r},t)=\frac{1}{\varepsilon_0}\rho(\vec{r},t) \quad  \quad \quad \quad \vec{\nabla}_r\cdot \vec{B}(\vec{r},t)=0 
\end{equation}
\begin{equation}\label{curl}
\vec{\nabla }_r\cross \vec{E}(\vec{r},t)=-\pdv{\vec{B}(\vec{r},t)}{t}  \quad  \quad \quad \quad \vec{\nabla }_r\cross \vec{B}(\vec{r},t)=\mu_0\qty(\vec{J}(\vec{r},t)+\varepsilon_0\pdv{\vec{E}(\vec{r},t)}{t})
\end{equation}
where we used the "vector-operator" $\vec{\nabla}_r:=(\pdv{}{x}, \pdv{}{y},\pdv{}{z})$, denoting by $\cdot$ the standard scalar vector product and by $\cross $ the standard cross product. $\varepsilon_0$ and $\mu_0$ are some universal constants fixed to satisfy experiments (the electric and magnetic permeabilities of free space, respectively). A separate law called the Lorentz force then rules the influence of these fields on the dynamics of the charge density and current.

By considering a region of space with no charges $\rho\equiv 0$ and no currents $\vec{J}=\vec{0}$, say in vacuum, we see that a time varying electric field must coexist with a time varying magnetic field and viceversa, suggesting a "self-sustained" wave that can travel in vacuum. In fact, we can take the curl $\vec{\nabla}_r$, of the equations in \eqref{curl} and using mathematical identities arrive to:
\begin{equation}
\mu_0\varepsilon_0\pdv[2]{\vec{E}(\vec{r},t)}{t}-\vec{\nabla}_r\cdot \vec{\nabla}_r\vec{E}(\vec{r},t)=0 \quad \quad \quad \mu_0\varepsilon_0\pdv[2]{\vec{B}(\vec{r},t)}{t}-\vec{\nabla}_r\cdot \vec{\nabla}_r\vec{B}(\vec{r},t)=0 
\end{equation}
which are standard wave equations for waves traveling at a speed $c:=\frac{1}{\sqrt{\varepsilon_0\mu_0}}$, known as the speed of light. This implies that an electric and magnetic field can indeed travel in free space as a so called electromagnetic wave. Then from the equations \eqref{div} in free space, we can further restrict this electromagnetic wave. We can easily verify that a traveling pair of waves $\vec{E}(\vec{r},t)=\vec{E}_0f(\vec{r}\cdot \hat{k}-ct)$, $\vec{B}(\vec{r},t)=\vec{B}_0f(\vec{r}\cdot \hat{k}-ct)$ with $\hat{k}$ a unit vector showing the direction of propagation of the wave and $\vec{E}_0$, $\vec{B}_0$ constants in space and time, do satisfy both wave equations for an arbitrary second differentiable function $f(s)$. Then, introducing them in \eqref{div}, we get that $\vec{k}\cdot\vec{E}f'(\vec{r}\cdot \hat{k}-ct)=0$ and $\hat{k}\cdot\vec{B}f'(\vec{r}\cdot \hat{k}-ct)=0$. Then, for a non-trivial $f$, it will need to happen that the electric and magnetic fields oscillate in the orthogonal plane to the propagation direction. Then using the first equation in $\eqref{curl}$ for this {\em ansatz}, we immediately get that $\hat{k}\cross \vec{E}_0 =c vec{B}_0$, meaning that the electric and magnetic fields will need to be orthogonal to each other and that $E_0=cB_0$. This then means that if we only describe the electric field for an electromagnetic wave of the given shape, we will have enough for a full description of the magnetic field. 

Arbitrary sinusoids moving in the spatial direction $\hat{k}:=\vec{k}/|\vec{k}|$ with an spatial angular frequency $|\hat{k}|$ and temporal angular frequency $w$, $f(\vec{r},t;\vec{k},w)=Acos(\vec{k}\cdot \vec{r}-wt)+Bsin(\vec{k}\cdot \vec{r}-wt)$ are convenient functions (as we will see), that can be shaped as the stated solution ansatz if we restrict $w(|\vec{k}|)=c|\vec{k}|$ (which is the so called dispersion relation). Their convenience comes mostly from the following succession of considerations:
\begin{itemize}
\item Sinusoids are a trivial shape a field generated by an oscillating source (like a dipole antenna irradiating the electromagnetic field) could be expected to have far away from their generation point.

\item Their behavior resembles the one of the classical light ray when describing how they refract into a medium, their interaction with classical optical elements etc., even if pure sinusoids are very unrealistic, since they occupy the whole space.

\item Since for any $A,B\in\R$, we can find real $R=\sqrt{A^2+B^2}$ and $\phi=tan^{-1}(B/A)$ such that $Acos(d)+Bsin(d)=Rcos(d+\phi)$, we see that an arbitrary sinusoid that solves the above equations can always be rewritten as a cosine:
\begin{equation}
f(\vec{r},t;\vec{k},w)=Acos(\vec{k}\cdot \vec{r}-wt)+Bsin(\vec{k}\cdot \vec{r}-wt)=Rcos(\vec{k}\cdot \vec{r}-wt+\phi)
\end{equation}
Now, since we have that such a sinusoid is equal to the real part of a plane wave with the same frequencies:
\begin{equation}
Rcos(\vec{k}\cdot \vec{r}-wt+\phi)=\mathcal{R}e\{R^{i\phi}e^{i(\vec{k}\cdot \vec{r}-wt)}\}
\end{equation}
It is typical in optics to talk about complex plane waves when describing an electromagnetic wave. Since we only add electric fields and multiply them by real scalars in Maxwell's equations, we can in general manipulate a complex electric field, just by remembering that in the end we only care about what is encoded by the real part (and discard the imaginary one). This turns out to be very advantageous when dealing with a plane wave, since for example, differentiation can be trivially done. Not only that, but as explained in Appendix $\alpha$, the complex plane wave is convenient when dealing with the polarization of the electric field (the evolution in time of the oscillation plane of this vector), since adding a phase shift between the vector components of the field, is reduced into multiplying a complex unit number to each component, turning basic discussions on polarization and optical elements into (complex) linear algebra.

\item It turns out that solving what happens to plane waves lets us know how any other acceptable initial conditions for the Maxwell Equations shall evolve. This is because, due to the Fourier decomposition theorem, an arbitrary real electric field $\vec{E}=(E_x,E_y,E_z)$ will be decomposable as a continuous series of sinusoids (or plane waves, if only the real part of such a decomposition is considered), for each component of the electric field $E_j$ with $j\in\{x,y,z\}$ will have its own Fourier series, for certain real functions $R_j,\phi_j$:
\begin{equation}
E_j(\vec{r},t)=\int_{\R^3} R_j(\vec{k},t)cos(\vec{k}\cdot\vec{r}+\phi_j(\vec{k},t))dk
=\mathfrak{R}e\qty{\int_{\R^3}R_j(\vec{k},t)^{i\phi_j(\vec{k},t)}e^{i\vec{k}\cdot\vec{r}}dk}
\end{equation}
Such that, for a certain complex vector field $\vec{E}_0(\vec{k},w)=(R_j(\vec{k},t)^{\phi_j(\vec{k},t)})_{j\in\{x,y,z\}}$, we will have that:
\begin{equation}
\vec{E}(\vec{r},t)=\mathfrak{R}e\qty{\int_{R^4}\vec{E}_0(\vec{k},t)e^{i\vec{k}\cdot\vec{r}}dk}
\end{equation}
Plugging it into the wave equation, one gets:
\begin{equation}
-k^2\vec{E}_0(\vec{k},t)-\frac{1}{c^2}\pdv[2]{\vec{E}_0}{t}=0
\end{equation}
whose solutions are $\vec{E}_0(\vec{k},t)=\vec{E}_0(\vec{k},t=0)e^{\pm iwt}$ with $w=|\vec{k}|c$. Then, we have that:
\begin{equation}
\vec{E}(\vec{r},t)=\mathfrak{R}e\qty{\int_{\R^3}\vec{E}_0(\vec{k},t=0)e^{i(\vec{k}\cdot\vec{r}-wt)}dk}
\end{equation}
Meaning that if we know the Fourier decomposition coefficients at the initial time for the initial electric field, we will be able to know all the subsequent time evolution of the arbitrary electromagnetic field. This is in fact a general qualitative consequence of the fact that Maxwell's equations are linear: if two functions are solution to the equation, then any linear combination of them will also be a solution, meaning if we have an initial condition that can be described as a linear combination (an uncountable one in this case), of initial conditions that we know how they evolve, we can get the time evolution of the initial conditions of interest by a "sum" of the decomposition states' evolution. 

\item By knowing the plane wave decomposition of the aribtrary initial conditions, not only we can know the evolution of the aribtrary initial electric field, but also the evolution of the magnetic field, since as we derived before, for each sinnusoid electric field, its coexisting magnetic field is directly determined in direction and magnitude, thus we can get it as well by a simple Fourier integral.

\end{itemize}

\mybox{{\bf Jones and Stokes Vectors to Describe the Polarization}\\

}

\section*{Anisotropic Medium for Light}
Now, it looks like this convenient description of light as an electromagnetic wave, where we only need to specify the electric field for sinusoidal plane waves will not be useful once we introduce other charges than the ones that generated the electromagnetic field. Yet, this is not really the case when we deal with linear dielectric materials, as most crystals are. It turns out, the propagation of light in a general crystal can be described as if it was traveling through an anisotropic "modified vacuum". 

Inside a material, electric fields cause a displacement on the electron clouds bound to each nuclei of its atoms, which generates tiny dipoles within each atom. All these tiny dipoles can statistically average to make macroscopically significant bound charge separations, which must be taken into account in macroscopic scale Maxwell equations. In fact, if the material has free electrons (like in metals), the effect will be even more dramatic. On the other hand, electrons bound to the neighbourhood of their nuclei have intrinsic angular momenta (due to spin and orbital angular momentum), which generate microscopic charge currents (magnetic dipoles), that can align and cause an average macroscopic bound current-like effect. Again, even more notorious will be the effect of free electrons. Then, inside a material, in general there will be two contributions to the total electric and magnetic fields: the one due to the bound charge- and current-densities $\rho_b$ and $\vec{J}_b$ and the one due to the rest of (free) charges and currents (including the external fields) $\rho_f$ and $\vec{J}_f$. By the linearity of Maxwell's equations, we can write the total fields as the sum of the fields generated by the bound and free charges: $\vec{E}=\vec{E}_b+\vec{E}_f$ and $\vec{B}=\vec{B}_b+\vec{B}_f$, such that we define the bound part of the electric field as the field generated by bound charge densities $\vec{\nabla}\cdot \vec{E}_b=\frac{1}{\varepsilon_0}\rho_b$ while the bound magnetic field as the one generated by the bound currents and the bound electric field: $\vec{\nabla}\cross \vec{B}_b=\vec{J}_b+\varepsilon_0\pdv{\vec{E}_b}{t}$. This expands Maxwell's equations as:
\begin{equation}\label{divBt}
\vec{\nabla}_r\cdot \vec{E_f}(\vec{r},t)=\frac{1}{\varepsilon_0}\rho_f(\vec{r},t) \quad  \quad\vec{\nabla}_r\cdot \vec{E_b}(\vec{r},t)=\frac{1}{\varepsilon_0}\rho_b(\vec{r},t) \quad \quad \vec{\nabla}_r\cdot \vec{B}_f(\vec{r},t)=0 \quad \quad \vec{\nabla}_r\cdot \vec{B}_b(\vec{r},t)=0 
\end{equation}
\begin{equation}\label{curlBt}
\vec{\nabla }_r\cross \vec{E}(\vec{r},t)=-\pdv{\vec{B}(\vec{r},t)}{t} 
\end{equation}
\begin{equation}
\vec{\nabla }_r\cross \vec{B}_f(\vec{r},t)=\mu_0\qty(\vec{J}_f(\vec{r},t)+\varepsilon_0\pdv{\vec{E}_f(\vec{r},t)}{t})\quad \quad  \vec{\nabla }_r\cross \vec{B}_b(\vec{r},t)=\mu_0\qty(\vec{J}_b(\vec{r},t)+\varepsilon_0\pdv{\vec{E}_b(\vec{r},t)}{t})
\end{equation}

In general, we are not interested in the explicit description of the bound charges. Instead, phenomenological constitutive equations are employed to describe their contribution for a given observed total electric and magnetic field $\vec{E}$, $\vec{B}$. For so called "linear" materials, these phenomenological macroscopic contributions are due to a linear transformation of the total fields: $\vec{E}_b=-\pmb{\chi_e} \vec{E}$ and $\pmb{\chi}\vec{B}_b= \vec{B}$ with $\pmb{\chi}:=\pmb{Id}+\pmb{\chi_m^{-1}}$, where $\pmb{\chi_e}$ and $\pmb{\chi_m}$ are material specific non-singular matrices called the electric and magnetic susceptibilities, $\pmb{(\cdot)^{-1}}$ denotes matrix inversion and $\pmb{Id}$ is the identity matrix\footnote{If the material has an homogeneous crystal structure, then these susceptibilities will be equal in all its inner points. However, we could also describe an inhomogeneous linear medium, by letting these susceptibilities depend on space (and time). }. 

Since the bound charges and currents manifest themselves as averaged electric and magnetic dipole moments, it is typical to describe their generated fields $\vec{E}_b$ and $\vec{B}_b$ in dipole moment density units (say, charge times displacement densities for the electric case), which can simply be done by multiplying with the permeability constants of free space. These are $\vec{P}:=-\varepsilon_0\vec{E}_b$ and $\vec{M}:=\frac{1}{\mu_0}\vec{B}_b$, called the polarization and magnetization fields. In order to unify the units it is typical to describe the non-bound charge fields as well in dipole moment density units, by defining $\vec{D}:=\varepsilon_0\vec{E}_f$ and $\vec{H}:=\frac{1}{\mu_0}\vec{B}_f$, which are called the displacement and magnetizing fields respectively.

Now, by knowing the phenomenological susceptibilities, we need not solve the equations for the total fields (all the expanded Maxwell equations). We have enough with the non-bound charge contributions (which do not require all the microscopic description as the bound charges in certain materials as we will see). For this, we just miss having the total fields in their function:
\begin{equation}\label{def1}
\vec{E}=\vec{E}_f+\vec{E}_b \Longrightarrow \varepsilon_0\vec{E}=\vec{D}+\varepsilon_0\vec{E}_b\Longrightarrow \vec{D}=\varepsilon_0(\pmb{Id}+\pmb{\chi_e})\vec{E} \Longrightarrow \vec{D}=\pmb{\varepsilon}\vec{E}
\end{equation} 
where $\pmb{\varepsilon}:=\varepsilon_0(\pmb{Id}+\pmb{\chi_e})$ is an invertible matrix, specific to the material relating the total and displacement fields known as the electric permeability of the material medium. Much the same way:
\begin{equation}\label{def2}
\vec{B}=\vec{B}_f+\vec{B}_b \Longrightarrow \frac{1}{\mu_0}\vec{B}=\vec{H}+\frac{1}{\mu_0}\vec{B}_b\Longrightarrow \vec{B}=\mu_0(\pmb{Id}+\pmb{\chi_m})\vec{H}\Longrightarrow \vec{B}=\pmb{\mu} \vec{H}
\end{equation} 
where $\pmb{\mu}:=\mu_0(\pmb{Id}+\pmb{\chi_m})$ is an invertible matrix specific to the material relating the total and magnetizing fields known as the magnetic permeability of the medium. Both permeabilities can be chosen to be complex if they are to model dephasing and medium absorption, yet, it can be proven that the real part must be diagonalizable.

With all this, we can now define the macroscopic Maxwell equations in a linear anisotropic medium as equations \eqref{def1} and \eqref{def2} together with:
\begin{equation}\label{div2}
\vec{\nabla}\cdot \vec{D}=\rho_f \quad \quad \quad \quad \quad \quad \vec{\nabla}\cdot \vec{H}=0
\end{equation}
\begin{equation}\label{curl2}
\vec{\nabla}\cross \vec{E}=-\pdv{\vec{B}}{t} \quad \quad \quad \quad \quad \quad \vec{\nabla}\cross \vec{H}=\vec{J}_f+\pdv{\vec{D}}{t}
\end{equation}
One of the reason why we do so, instead of using the total field equations, is because we see that this version of Maxwell's equations have a similar look to the equations in free space. In particular, if for instance we consider a linear and homogeneous material where the susceptibilities are constant scalars within the material (the three eigenvalues are the same and are constant), what we call an isotropic material, we would see that the first equation in \eqref{curl2} would become in $\vec{\nabla}\cross\vec{D}=-\varepsilon\mu\pdv{\vec{H}}{t}$, which would now leave a system exactly the same as the microscopic equations (plus the consitutive phenomenological equations to get the full fields). In such a case, we see that if we consider a dielectric material: meaning that there is no free charge that could carry electric current (most crystals used in optics), inside the material, the free charge density and current will be null (only bound charges and currents will be present). Thus, by taking the curl of the equations in \eqref{curl2} as we did in free space, we would arrive to two equations:
\begin{equation}
\varepsilon\mu\pdv[2]{\vec{D}(\vec{r},t)}{t}-\vec{\nabla}_r\cdot \vec{\nabla}_r\vec{D}(\vec{r},t)=\vec{0} \quad \quad \quad \varepsilon\mu\pdv[2]{\vec{H}(\vec{r},t)}{t}-\vec{\nabla}_r\cdot \vec{\nabla}_r\vec{H}(\vec{r},t)=\vec{0}
\end{equation}
where we recognize vector wave equations where the propagation of the waves is now at a speed $\frac{1}{\sqrt{eig(\varepsilon\mu)_j}}$ instead of $c$. For the rest of matters, such an isotropic material would behave as free space, with the same consequences for the description of self-sustained electromagnetic waves in terms of $e^{\vec{k}\cdot\vec{r}-wt}$ plane-waves.

However, we know that most crystals show a different light propagation as a function of the propagation direction. This is because, their susceptibilities are not scalars. These are called anisotropic materials. In these materials we are no longer able to obtain a typical wave equation (since curl and matrix multiplication do not commute), but we can still see that for dielectric materials (no free charges nor currents), plane wave solutions that satisfy our intuitions for light in free space still exist. In particular, let us assume solutions $\vec{D}=\vec{D}_0e^{\vec{k}\cdot\vec{r}-wt}$ and $\vec{H}=\vec{H}_0e^{\vec{k}\cdot\vec{r}-wt}$ with $\vec{D}_0,\vec{H}_0\in\C^3$ arbitrary complex constants to account for the directions of oscillation and the phase difference between components. Note we leave the speed of propagation of the equal phase planes $v=w/|\vec{k}|$ as a free parameter that will be fixed by the equations later. We first observe that since $\vec{E}=\pmb{\varepsilon^{-1}}\vec{D}$ and $\vec{B}=\mu\vec{H}$, if $\vec{D}$ and $\vec{H}$ are plane waves, $\vec{E}=e^{\vec{k}\cdot\vec{r}-wt}\pmb{\varepsilon^{-1}}\vec{D}_0$ and $\vec{B}=e^{\vec{k}\cdot\vec{r}-wt}\pmb{\mu}\vec{H}_0$ will as well be plane waves with the same frequencies, even if they will not satisfy all the intuitions (which is actually part of the reason why we shaped the macroscopic equations for dielectrics this way). In particular we will have by plugging $\vec{D}$ and $\vec{H}$ in \eqref{div2}, that $\vec{k}\cdot\vec{D}=0$: the propagation direction $\vec{k}$ will need to be orthogonal to the oscillation direction $\vec{D}_0$. The same goes for $\vec{H}$ (and actually also for $\vec{B}$, yet not for $\vec{E}$). On the other hand, plugging them in \eqref{curl2}, we find:
\begin{equation}\label{Fcurl}
\vec{k}\cross\vec{E}=w\vec{B} \quad \text{           and           }\quad  \vec{k}\cross\vec{H}=-w\vec{D}
\end{equation}  
These mean that $\vec{E}$ and $\vec{B}$ are orthogonal to each other, $\vec{H}$ and $\vec{D}$ as well (which are also orthogonal to the propagation direction as light in free space), and $|\vec{H}_0|=v|\vec{D}_0|$ and $|\vec{E}_0|=v|\vec{B}_0|\Leftrightarrow |\pmb{\varepsilon^{-1}}\vec{D}_0|=v|\pmb{\mu}\vec{H}_0|$ (which for isotropic media would yield $v^2=\frac{1}{\varepsilon\mu}$ as expected), leaving the "magnetic part" completely determined by the "electric part" of the plane waves. 

\section*{The Dispersion Relation in Anisotropic Dielectric Media}
We just miss the question of how $w$ and $\vec{k}$ of the plane wave solutions are mutually restricted (the dispersion relation), or equivalently, which is the phase propagation speed $v=w/|\vec{k}|$ for a plane wave with temporal frequency $w$. The ratio of this speed relative to free space $n:=c/v=c|\vec{k}|/w$ is called the refractive index. In answering this question, we will finally arrive to the Conical Refraction phenomenon.

For this, retake equations \eqref{Fcurl} and plug in the constitutive equations \eqref{def1} and \eqref{def2} to get:
\begin{equation}\label{honek}
\vec{k}\cross \vec{E}=w\pmb{\mu}\vec{H} \quad \quad \quad \vec{k}\cross\vec{H}=-w\pmb{\varepsilon}\vec{E}
\end{equation} 
If we now define the cross product operator:
\begin{equation}
\pmb{k}:=\begin{pmatrix}
0& -k_z& k_y\\
k_z & 0 & -k_x\\
-k_y & k_x & 0
\end{pmatrix}
\end{equation}
such that $\vec{k}\cross\vec{a}=\pmb{k}\vec{a}$ for any $\vec{a}\in\C^3$.  
Using the inverses of $\pmb{\varepsilon}$,$\pmb{\mu}$ and applying $\pmb{k}$ in each side of equations \eqref{honek}, we can easily rewrite them as:
\begin{equation}
\qty[\pmb{k}\pmb{\varepsilon^{-1}}\pmb{k}+w^2\pmb{\mu}]\vec{H}=\vec{0} \quad \quad \quad \qty[\pmb{k}\pmb{\mu^{-1}}\pmb{k}+w^2\pmb{\varepsilon}]\vec{E}=\vec{0}
\end{equation}
These are the so called complex vector wave equations for anisotropic media. For a non trivial solution ("not all zero") to exist, the determinant of the matrices preceding $\vec{H}$ and $\vec{E}$ must be zero:
\begin{equation}\label{wave2}
det\qty(\pmb{k}\pmb{\varepsilon^{-1}}\pmb{k}+w^2\pmb{\mu})=0 \quad \quad\quad det\qty(\pmb{k}\pmb{\mu^{-1}}\pmb{k}+w^2\pmb{\varepsilon})=0
\end{equation}
Since the matrices $\pmb{\mu}$ and $\pmb{\varepsilon}$ are known, these equations will give the restriction relating the allowed plane-wave $vec{k}$ and $w$ for this medium, known as the dispersion relation $w(\vec{k})$. In reality they are equivalent conditions, since both $\vec{H}$ and $\vec{E}$ must coexist, so we can choose the simplest one of them as a function of how difficult inverting the $\pmb{\mu}$ and $\pmb{\varepsilon}$ are.

In our case, we will choose to study crystals that barely have magnetic response, which are the typical ones employed in optical elements. This will mean that $\pmb{\mu}\simeq \mu_0\pmb{Id}$. As such, the simplest equation to develop the dispersion relation will be the second one in \eqref{wave2}. Let us choose the reference frame in which the coordinate axes are aligned with the eigenvectors of $\pmb{\varepsilon}$ (which can be orthogonal, for its hermiticity). By defining the ordered eigenvalues of $\pmb{\varepsilon}$ as $\varepsilon_x, \varepsilon_y, \varepsilon_z$, it is easy to get from the second equation of \eqref{wave2} with $\vec{k}=(k_x,k_y,k_z)$:
\begin{equation}
det\begin{pmatrix}
w^2\mu_0\varepsilon_x-(k_y^2+k_z^2) & k_xk_y & k_xk_z\\
k_yk_x & w^2\mu_0\varepsilon_y-(k_z^2+k_x^2) & k_yk_z\\
k_xk_z & k_yk_z & w^2\mu_0\varepsilon_z-(k_x^2+k_y^2)
\end{pmatrix}=0
\end{equation}
which after developing the equation, dividing by $w^2$ and by $\mu_0$ both sides and multiplying the whole by $c^6\mu_0$, leads to the next equation if grouping by $w$:
\begin{equation}\label{Fres}
w^4\varepsilon_x\varepsilon_y\varepsilon_z-w^2c^2\qty[k_x^2\varepsilon_x(\varepsilon_y+\varepsilon_z)+k_y^2\varepsilon_y(\varepsilon_x+\varepsilon_z)+k_z^2\varepsilon_z(\varepsilon_x+\varepsilon_y) ]+|\vec{k}|^2c^4\qty[\varepsilon_xk_x^2+\varepsilon_yk_y^2+\varepsilon_zk_z^2]=0
\end{equation}
For a given constant $w$, this is a fourth order surface in the possible wavevector $\vec{k}=(k_x,k_y,k_z)$ space, known as Fresnel's surface. It can readily be seen that by fixing a certain $w$, and given the direction of $\vec{k}$, just leaving its magnitude $|\vec{k}|^2$ as unknown, equation \eqref{Fres} is a quadratic equation with real coefficients. This means, in general there will be two possible spatial frequencies $|\vec{k}|$, thus two possible phase velocities $v=w/|\vec{k}|$ and two possible refractive indices $n=c/v$ for each given direction in $k$-space. This would be opposed to the isotropic or free space plane waves, where there was only a single acceptable wavelength and refractive index. In fact, this is why for material media where at least one of the eigenvalues of $\pmb{\varepsilon}$ is different to the rest, if a single light ray incides the crystal, it will propagate inside it as two light rays with different refractive indices. This phenomenon is known as birrefringence. Yet, it turns out Fresnel's surface will have some singular directions where there is only a single allowed refractive index, which will turn out to yield the prediction of Conical Refraction. For this however, we miss a last little step, by which we will realize that these constant $w$ surfaces for the function $w(\vec{k})$ implicitly defined in \eqref{Fres}, also tell us the direction in which an actual light ray (composed of many plane waves) will propagate.

\mybox{{\bf Group Velocity}\\
Plane waves are very unrealistic since they have an infinite extent in space and time. Instead, we could write a more general finite extent field following the general Fourier integral as a (continuous) superposition of plane waves:
\begin{equation}
\vec{H}(\vec{r},t)= \frac{1}{(2\pi)^4}\int_{\R^4}\vec{H}(\vec{k},w)e^{i(\vec{k}\cdot\vec{r}-wt)}\; dw\;dk_x\;dk_y\;dk_z
\end{equation}
We already saw that not all $w$ and $\vec{k}$ pairs are allowed, which had to follow $w(\vec{k})$ implicitly defined in \eqref{Fres}. Thus in reality:
\begin{equation}
\vec{H}(\vec{r},t)= \frac{1}{(2\pi)^3}\int_{\R^3}\vec{H}(\vec{k})e^{i(\vec{k}\cdot\vec{r}-wt)} \;dk_x\;dk_y\;dk_z
\end{equation}
Now, a typical light ray like a laser beam is very monochromatic, it has a sharply defined wavelength $\lambda_c$, and has a defined phase propagation direction, implying there is a dominant plane wave in its decomposition of wavenumber $\vec{k}_c$ of magnitude $|\vec{k}_c|=\frac{2\pi}{\lambda_c}$. If so, the Fourier coefficients $\vec{H}(\vec{k})$ will be sharply peaked around $\vec{k}_c$, making only $\vec{k}$ around $\vec{k}_c$ relevant. This implies, we could consider only the first two terms of the Taylor series of $w(\vec{k})$ around $\vec{k}_c$:
\begin{equation}
w(\vec{k})\simeq w(\vec{k}_c)+\vec{\nabla}_kw(\vec{k})\Big\rvert_{\vec{k}_c}\cdot\qty(\vec{k}-\vec{k}_c)+\cdots
\end{equation}
By writing $\vec{k}=\vec{k}_c+(\vec{k}-\vec{k}_c)$, we can rewrite the Fourier expansion as:
\begin{equation}
\vec{H}(\vec{r},t)=e^{i(\vec{k}_c\cdot\vec{r}-w(\vec{k_0})t)}\int_{\R^3}\vec{H}(\vec{k})\;e^{i\qty((\vec{k}-\vec{k}_c)\cdot [\vec{r}-\vec{\nabla}_kw(\vec{k}_c)t])}\;dk_x\;dk_y\;dk_z
\end{equation}
The integral varies very slowly in space compared with the factored out plane wave, which has a very high frequency in space $|\vec{k}_c|$ compared with the relevant plane waves for the integral, since we assumed$|\vec{k}-\vec{k}_c|<<|\vec{k}_0|$. Then we realize that the factored out plane wave is a fast sinusoid that is modulated by the slowly varying envelope integral. As such, in order to know the direction of motion of the light packet, we will need to see the direction in which the envelope is moving. We can fix our attention in a particular point of the envelope by tracking the point in space-time which has the same integral, meaning the same value for $\vec{r}-\vec{\nabla}_kw(\vec{k}_c)t=c$ for a certain constant $c$. Finally, we will get the velocity vector for the displacement of the fixed point of the envelope by a time derivative:
\begin{equation}
\vec{v}_g(\vec{k}_c)=\vec{\nabla}_kw(\vec{k})\Big\rvert_{\vec{k}_c}
\end{equation}
We call $\vec{v}_g(\vec{k})$ the group velocity of the light ray of main frequency $\vec{k}$
}

Thus, a light ray of main frequency $\vec{k}$ will propagate in space in the direction of the gradient of $w(\vec{k})$ at that point, which geometrically implies it will move following the orthogonal vector to the tangent plane of the constant $w$ Fresnel surface level of $w(\vec{k})$. Note interestingly that a plane wave may have a phase velocity vector (pointing in the direction in which equal phase planes or wavefronts flows) $\vec{v}(\vec{k})=w/|\vec{k}|^2\vec{k}$ that is different to the actual propagation direction of its composing beam $\vec{v}_g(\vec{k})=\vec{\nabla}_kw(\vec{k})$, both in magnitude and direction\footnote{It can be proven, following the intuition for a moving wave-packet that the direction of energy flow matches the group velocity's direction.}. This "sideways-walking" rays are called extraordinary rays (for they will not obey exactly Snell's law).



\section*{Conical Refraction}

There are three types of crystals as a function of their electric permitivity $\pmb{\varepsilon}$. If the three eigenvalues $\varepsilon_x=\varepsilon_y=\varepsilon_z$ are equal, the medium is isotropic and the crystal will have a cubic lattice. If there is one eigenvalue that is different to the other two $\varepsilon_{xy}:=\varepsilon_x=\varepsilon_y\neq \varepsilon_z$, then the crystal is called uniaxial for the reason we will see in a moment. These include rhombohedral, tetrahedral and hexagonal lattices, which have a main rotational symmetry axis which will be the eigenvector of $\varepsilon_z$, called in this case, the optical axis of the crystal. By setting these $\varepsilon_j$ in Fresnel's equation \eqref{Fres}, it factors as: $(k^2c^2-w^2\varepsilon_{xy})\qty(c^2(\varepsilon_zk_z^2+\varepsilon_{xy}(k_x^2+k_y^2))-w^2\varepsilon_z\varepsilon_{xy})=0$, which is solved by the two quadratic equations:
\begin{equation}
\frac{k_x^2+k_y^2+k_z^2}{\varepsilon_{xy}}=\frac{w^2}{c^2} \quad \quad \quad \frac{k_z^2}{\varepsilon_{xy}}+\frac{k_x^2+k_y^2}{\varepsilon_z}=\frac{w^2}{c^2}
\end{equation}
which are the equation of a sphere and an ellipsoid, where one of the semi-axes of the ellipsoid matches the radious of the sphere. This is depicted for a fixed $w$ in Figure \ref{fig:concentric}.a. This means that there is one axis in the crystal (the eigenspace for $\varepsilon_z$), along which for a fixed $w$, there is only one allowed $\vec{k}$, but for all the rest of directions, two $\vec{k}$ happen. This causes that when a single ray enters the crystal, unless it was directed along the optical axis, it will be refracted with two different phase velocities. In fact, since what in an interface between media what must be continuous is the tangent component of the wave-vector to the interface, this means that the ray will be doubly refracted with different directions due to the fact that this restriction will make the orthogonal to the tangent planes of the ellipsoid and the circle be different. Thus, these crystals are called uniaxial because they have a single axis along which no birefringence occurs.

It was found in the 17th century that there were some other crystals where birrefringence was also present, but where there were two axes along which birrefringence did not occur. These were called biaxial crystals, and are due to those lattices in which there is no clear symmetry axis: triclinic, monoclinic and rhombic systems. What was not still clear in that century is that along these two axes there is no normal refraction either (as happens in the optical axis of uniaxial crystals). Instead, it is there where conical refraction takes place.

In order to plot its Fresnel surface, we can first obtain the shape of its intersection with the coordinate planes (the axes aligned with the eigenvectors of $\pmb{\varepsilon}$). we will assume $\varepsilon_x<\varepsilon_y<\varepsilon_z$. If we check the $xy$ plane by setting $k_z=0$, we get $(k^2c^2-\varepsilon_zw^2)(c^2(\varepsilon_xk_x^2+\varepsilon_yk_y^2)-w^2\varepsilon_x\varepsilon_y)=0$, which are the circumference $k_x^2+k_y^2=\varepsilon_zw^2/c^2$ and the ellipse $\frac{k_x^2}{\varepsilon_y}+\frac{k_y^2}{\varepsilon_x}=w^2/c^2$, which do not intersect for $\varepsilon_z$ is the biggest eigenvalue of $\pmb{\varepsilon}$. The intersection with the $yz$ plane also gives a circumference and an ellipse that do not intersect (the circumference has a radius $\varepsilon_x$, the smallest eigenvalue). The intersection with the $xz$ plane on the other hand, gives an ellipse and a circumference that do intersect, since the circumference now has the intermediate radius $\varepsilon_y$ between the two semi-axes $\varepsilon_x$ and $\varepsilon_y$ of the ellipse. You can see this in Figure \ref{fig:concentric}.b. For continuity and symmetry we can imagine the rest of the surface replicating this picture. It is a figure with four singular points at which the two constant $w$ sheets touch, being two of them along one same axis and the other two along another same axis. These are the two axes along which plane waves can propagate only with a single phase velocity (single refractive index), thus the name biaxial crystal. Note how around each singular point each of the constant $w$ sheets forms a cone, the vertices of which meet in the singular point, leaving the singular point in the center of a "diabolo"-like shape.

Now, we said that the normal to the constant $w$ surface indicated for a given wavevector the propagation direction of a nearly monochromatic light ray with such a main plain wave in its decompostion. It turns out however the tangent plane of the "diabolo" point is not well defined, thus neither its normal vector. However, since a real light ray, say a Gaussian laser beam, has other wavevector components around the singularity wavevector, we can predict by continuity which will be the propagation direction such a light ray with singular wavevector. For this, note that the closer we get to the singularity following one of the sheets, the closer will be the ensemble of the tangent planes to a perfect cone, in the sense that the ensemble of normal vectors to these tangent planes around the singularity will each time be closer to a conical bouquet of group velocities. Thus, a light beam incident on a biaxial crystal cut orthogonally to one of its optical axes, a beam with its main wave-vector aligned with the optical axis, will refract not as a single beam, nor as two beams (as in birrefringence), but as a cone of "infinite" rays. Thus, inside the crystal the light beam will open into a hollow cone of light. Once the cone of light arrives to the end face of the crystal, it will get out as a hollow ring of light. This is what Hamilton called Conical Refraction.

\mybox{{\bf The Relevance of Conical Refraction in the History of Physics}\vspace{0.2cm}\\
Back in the late 18-th century, a vivid battle was taking place in physics, between those who defended light was an ensemble of discrete corpuscles and those who defended it was a wave. The former had the authority of Sir Isaac Newton backing them, while the latter really got momentum when Christiaan Huygens developed his theory of wavefronts in 1690, which allowed the description of birefringence for uniaxial crystals and how there was an axis along which this did not happen, a phenomenon that had no previous explanation. Yet, David Brewster found in 1813 that there were some crystals where there were two axes in which birefringence did not happen (biaxial crystals), which Huygens' theory could not explain. Moreover, Pierre Simon Laplace and Etienne Louis Malus achieved in the description of birefringence in terms of corpuscles. In 1801, Thomas Young revealed the light interference and slit difraction phenomenon, which further suggested the wave nature of light. Finally, in 1819, Agustin-Jean Fresnel presented his theory of light, by which he was able to explain birefringence in both uniaxial and biaxial crystals together with interference and diffraction. For this he derived the equation we know today by his name and which we derived from Maxwell's equation (equation \eqref{Fres}), just from geometrical intuitions, when Maxwell was still about to be born. His theory was yet very involved and still the balance needed a further push to decide the winner posture. \\

This last push came from William Rowan Hamilton, who deeply studied Fresnel's surfaces and in 1832 found what we have seen: the singularities of biaxial crystal surfaces had a double cone apex shape, which meant a proper light beam crossing one of the optical axes of a biaxial crystal should refract as a cone and get out of the crystal as a hollow cylinder. A few weeks later, after being asked by Hamilton, Humphrey Lloyd was able to experimentally verify this fact, making Hamilton's mathematical discovery, one of the first examples in the history of science where a mathematical prediction preceded its experimental discovery. This was a hit in the struggle in favor of the wave description of light. Amazingly, after this, conical refraction remained as a pure curiosity and lost the interest of the community for about a century, with only sporadic contributions by scientists like Johann Christian Poggendorff (in 1839) or Chandrasekhara Venkata Raman (in 1941) among others, who discovered that in fact the ring was split in two light rings by a dark ring, known today as the Poggendorff dark ring, or the fact that it was in the focal plane where the ring was most clearly visible, such that the ring contracts until a light point as we get away from itthe focal plane, achieving a maximum axial intensity in which today is known as the Raman spot.\\ 

The disappearance of CR was in part because a full diffraction theory was to be developed in order to account for the complexity of the phenomenon, and in part because the experimental fineness for a decent account of its verification would only be achievable in the last half of the 20-th century, with the advance of laser and crystal engineering technologies. It was in this last period of the century that a revival of the interest in CR emerged, when the complete mathematical formulation for the phenomeon (as we will review in a moment) and uncountable applications for it (among others the polarimeter developed in this thesis) were suggested.
}

%$(k^2c^2/w^2-\varepsilon_{xy})\qty(c^2/w^2(\varepsilon_zk_z^2+\varepsilon_{xy}(k_x^2+k_y^2))-\varepsilon_z\varepsilon_{xy})=0$ 

\section*{A.2. The Polarization of each Ray along the Cone:\\ The Belsky-Khapalyuk-Berry Description of the Phenomenon }
\addcontentsline{toc}{subsection}{2. Berry's Mathematical Description of the Phenomenon }
En más o menos detalle guredozulez azaldu fenomenoan matematikie eta batezbe heldu formula finalatara

\section*{A.3. Simulating the Phenomenon }
\addcontentsline{toc}{subsection}{3. Simulating the Phenomenon }
Azaldu GPU/CPU tradeoffa, zelan implemente doten eta jarri imagenak. Azaldu zelan si lienar pol tal, si cricular pol tal imagenakaz. Sugeridu zer alko genun ein orduen linear polrztion aldaketak antzemateko.

\section*{A.4. A Natural Polarimeter}
\addcontentsline{toc}{subsection}{4. A Natural Polarimeter }
Esan zelan alko zendun argixen polarizaziño tal danak deskribatu einde bi besogaz et al, baia zelan simplifike al dan ze kiralidade temak LPgaz nahiko eta hori da polarimetroan merkatu handixetako bat.

\section*{Apppendix $\alpha$: Jones and Stokes Representations of Polarization}
\addcontentsline{toc}{subsection}{Apppendix $\alpha$: Jones and Stokes Representations of Polarization }



\newpage

\fancyhead[OL]{\em Part A: Conical Refraction Essentials}
\fancyhead[OR]{ Section title}

\fancyhead[EL]{ Section title}
\fancyhead[ER]{\em Bachelor's Thesis: Conical Refraction Polarimeter}
Ba hori



%%%%%%%%%%%%%%%%%%%%%%%%%%%%%%%%%%%%%%%%%%%%%%%%%%%%%%%%%%%%%%%%%%%%%%%%%%%%%%%%%%%%%%%%%%%%%%%%%%%%%%%%%%%%%%%%%%%%%%%%%%%%%%%%%%%%%%%%%%%%%%%%%%%%%%%%%%%%%%%%%%%%%%%%%%%%%
\newpage


%\begin{kapituloBerria}{Part B}{ The Path Towards a Fully Functional Portable Prototype}
%Designing the Hardware
%\end{kapituloBerria}
%\newpage
%\fancyhead[L]{\null}
%\fancyhead[R]{\null}
%\null
%\clearpage


\section*{\centering \huge{Part B: Designing the Hardware}\vspace{-0.3cm}}
\noindent\rule{\textwidth}{0.4pt}

\addcontentsline{toc}{section}{Part B: Designing the Hardware}

\section*{B.1. Experimental Implementations of the Polarimeter}
\addcontentsline{toc}{subsection}{1. Experimental Implementations of the Polarimeter}
Azaldu bakjoitza dibujo bategaz eta identifike problemak, edo gubazun hori hobeto alrgazki batzuk ipiñitze hurrengo seksiñoan, esaten posibles razones a las ke se deba, al dozu citeu ah artikuloa ta beran imagena bebai.


\section*{B.2. Main Non-Idealities in the Experimental Rings}
\addcontentsline{toc}{subsection}{2. Main Non-Idealities in the Experimental Rings}

\section*{B.3. Final Proofs of Concept}
\addcontentsline{toc}{subsection}{3. Proofs of Concept}
Bat ya einde egioin zana, en el ke implementé la GUI akella, argazki bat, con fotos de la GUI etc.

Bestie, lortzen badogun eitzie housiñegaz y tal, planoa ein eskuz ni ke sea

COST ANALYSIS EGIN

\section*{Appendix $\beta$: Employed Optical Elements}
\addcontentsline{toc}{subsection}{Appendix $\beta$: Employed Optical Elements}



\newpage


\fancyhead[OL]{\em Part B: Designing the Hardware}
\fancyhead[OR]{ Section title}

\fancyhead[EL]{\em Section title}
\fancyhead[ER]{\em Bachelor's Thesis: Conical Refraction Polarimeter}
Ba hori



%%%%%%%%%%%%%%%%%%%%%%%%%%%%%%%%%%%%%%%%%%%%%%%%%%%%%%%%%%%%%%%%%%%%%%%%%%%%%%%%%%%%%%%%%%%%%%%%%%%%%%%%%%%%%%%%%%%%%%%%%%%%%%%%%%%%%%%%%%%%%%%%%%%%%%%%%%%%%%%%%%%%%%%%%%%%%
\newpage


%\begin{kapituloBerria}{Part C}{ From Crescent Ring Images to Polarization Angle Differences }
%Designing the Software
%\end{kapituloBerria}
%\newpage
%\fancyhead[L]{\null}
%\fancyhead[R]{\null}
%\null
%\clearpage


\section*{\centering \huge{Part C: Designing the Software}\vspace{-0.3cm}}
\noindent\rule{\textwidth}{0.4pt}
\newpage

\addcontentsline{toc}{section}{Part C: Designing the Software}

Objetiboak argitu, ta komplikaziñoak 

\section*{C.1. Artificial Noise Generation}
\addcontentsline{toc}{subsection}{1. Artificial Noise Generation}

\section*{C.2. Simulated Image Datasets }
\addcontentsline{toc}{subsection}{2. Simulated Image Datasets}

\section*{C.2. Preprocessing }
\addcontentsline{toc}{subsection}{2. Preprocessing}

\section*{C.3. Embedding Space Algorithms}
\addcontentsline{toc}{subsection}{3. Embedding Space Algorithms}

\subsection*{C.3.1. Data Manifold Dimension Identification}
\addcontentsline{toc}{subsubsection}{3.1. Data Manifold Dimension Identification}

\subsection*{C.3.2. Metric Learning }
\addcontentsline{toc}{subsubsection}{3.2. Metric Learning}

\subsection*{C.3.3. Nearest Neighbors }
\addcontentsline{toc}{subsubsection}{3.3. Nearest Neighbors}

\section*{C.3. Geometric Algorithms}
\addcontentsline{toc}{subsection}{3. Geometric Algorithms}

Tos los geometricos y los optimizadores implementados.

\section*{C.4. Simulation Fitting Algorithms}
\addcontentsline{toc}{subsection}{4. Simulation Fitting Algorithms }

\section*{C.5. The Best Algorithm and Preprocessing}
\addcontentsline{toc}{subsection}{5. Looking for the Best Algorithm and Pre-processing }

\section*{Appendix $\gamma$: The Implemented Optimizers}
\addcontentsline{toc}{subsection}{Appendix $\gamma$: The Implemented Optimizers }

\fancyhead[OL]{\em Part C: Designing the Software}
\fancyhead[OR]{ Section title}

\fancyhead[EL]{\em Section title}
\fancyhead[ER]{\em Bachelor's Thesis: Conical Refraction Polarimeter}
Ba hori


%%%%%%%%%%%%%%%%%%%%%%%%%%%%%%%%%%%%%%%%%%%%%%%%%%%%%%%%%%%%%%%%%%%%%%%%%%%%%%%%%%%%%%%%%%%%%%%%%%%%%%%%%%%%%%%%%%%%%%%%%%%%%%%%%%%%%%%%%%%%%%%%%%%%%%%%%%%%%%%%%%%%%%%%%%%%%
\newpage

%
%\begin{kapituloBerria}{Part D}{ Performance Analysis, Commercial Polarimeters and Niche Outlining  }
%The Final Device
%\end{kapituloBerria}
%\newpage
%\fancyhead[L]{\null}
%\fancyhead[R]{\null}
%\null
%\clearpage


\section*{\centering \huge{Part D: The Final Device}\vspace{-0.3cm}}
\noindent\rule{\textwidth}{0.4pt}

\addcontentsline{toc}{section}{Part D: The Final Device}

\section*{D.1. Performance on Experimental Data }
\addcontentsline{toc}{subsection}{1. Performance on Experimental Data }

\section*{D.2. Commercial Polarimeters }
\addcontentsline{toc}{subsection}{2. Commercial Polarimeters }

\section*{D.3. Potential Niches }
\addcontentsline{toc}{subsection}{3. Potential Niches }

\newpage

\fancyhead[OL]{\em Part D: The Final Device}
\fancyhead[OR]{ Section title}

\fancyhead[EL]{\em Section title}
\fancyhead[ER]{\em Bachelor's Thesis: Conical Refraction Polarimeter}

\newpage
\fancyhead[OL]{}
\fancyhead[OR]{}
\fancyhead[EL]{}
\fancyhead[ER]{}

\section*{\centering \huge{Conclusions}\vspace{-0.3cm}}
\noindent\rule{\textwidth}{0.4pt}
\addcontentsline{toc}{section}{Conclusions}



\newpage
\begin{thebibliography}{1}
\addcontentsline{toc}{section}{References}


\bibitem{where}

\end{thebibliography}


\end{document}
